% Part: computability
% Chapter: recursive-functions
% Section: composition

\documentclass[../../../include/open-logic-section]{subfiles}

\begin{document}

\olfileid{cmp}{rec}{com}
\olsection{સંયોજન}

પ્રાકૃતિક સંખ્યાઓનાં એક-આર્ગ્યુમેન્ટવાળાં બે વિધેયો $f$ અને $g$ હોય,
તો તેમનું સંયોજન કરી શકાય: $h(x) = f(g(x))$. આમ નવું વિધેય~$h(x)$,
$f$ અને~$g$ના \emph{સંયોજન}થી વ્યાખ્યાયિત થાય છે. હવે આ રીતને
એકથી વધુ આર્ગ્યુમેન્ટ ધરાવતાં વિધેયો માટે સામાન્ય બનાવવા માગીએ છીએ.

તેની એક રીત આ છે: ધારો કે $f$ને $k$ આર્ગ્યુમેન્ટ છે અને
$g_0$, \dots, $g_{k-1}$ એવાં $k$ વિધેયો છે, જેમાંના દરેકને
$n$ આર્ગ્યુમેન્ટ છે. તો નવું $n$-આર્ગ્યુમેન્ટવાળું વિધેય $h$ આ રીતે
વ્યાખ્યાયિત કરી શકાય:
\[
h(x_0, \dots, x_{n-1}) =
f(g_0(x_0, \dots, x_{n-1}), \dots, g_{k-1}(x_0, \dots, x_{n-1}))
\]
$f$ અને દરેક $g_i$ સંગણનીય હોય તો $h$ પણ સંગણનીય છે:
$h(x_0, \dots, x_{n-1})$ ગણવા માટે પહેલાં દરેક $i = 0$,
\dots,~$k-1$ માટે $y_i = g_i(x_0, \dots, x_{n-1})$ ગણો.
પછી આ મૂલ્યો $f$માં આપીને
$h(x_0, \dots, x_{n-1}) = f(y_0, \dots, y_{k-1})$ ગણો.

\sourcecorrection{OLCMP-002}{સ્થિર અંગ્રેજી મૂળના આ વાક્યમાં
$h(x_0,\dots,x_{k-1})$ લખાયું છે. અહીં $h$ને $n$ આર્ગ્યુમેન્ટ છે;
અતઃ ડાબી બાજુનો અંતિમ સૂચક $n-1$ હોવો જોઈએ. ગુજરાતી લખાણમાં
આ જ સૂત્ર $h(x_0,\dots,x_{n-1})=f(y_0,\dots,y_{k-1})$ રાખ્યું છે.}

પહેલેથી મળેલાં વિધેયો વાપરીને નવું વિધેય ગણીએ ત્યારે શું થાય છે તેનું
આ વર્ણન કદાચ વધારે પડતું મર્યાદિત લાગે. એક કારણ એ છે કે ક્યારેક
આપણે વિધેયના બધા આર્ગ્યુમેન્ટ વાપરતા નથી: $\Add$ની આદિમ
પુનરાવર્તી વ્યાખ્યા માટે $g(x, y, z) = \Succ(z)$ આપ્યું ત્યારે પણ
એવું જ થયું હતું. ધારો કે નીચેના વિધેયો વાપરવાની છૂટ છે:
\[
\Proj{n}{i}(x_0, \dots, x_{n-1}) = x_i
\]
$\Proj{k}{i}$ને \emph{પ્રક્ષેપ} વિધેયો કહે છે:
$\Proj{n}{i}$ને $n$ આર્ગ્યુમેન્ટ છે. હવે $g$ને આ રીતે
વ્યાખ્યાયિત કરી શકાય:
\[
g(x, y, z) = \Succ(\Proj{3}{2}(x, y, z)).
\]
અહીં $f$ની ભૂમિકા $1$-આર્ગ્યુમેન્ટવાળું વિધેય $\Succ$ ભજવે છે,
તેથી $k=1$. અને $g_0$ની ભૂમિકા એક $3$-આર્ગ્યુમેન્ટવાળું વિધેય
$\Proj{3}{2}$ ભજવે છે. પરિણામે મળતું $3$-આર્ગ્યુમેન્ટવાળું વિધેય
તેના ત્રીજા આર્ગ્યુમેન્ટનો ઉત્તરગામી આપે છે.

પ્રક્ષેપ વિધેયો આર્ગ્યુમેન્ટોનો ક્રમ બદલવા અથવા બે આર્ગ્યુમેન્ટોને
એક જ આર્ગ્યુમેન્ટથી ભરવા વડે પણ નવાં વિધેયો વ્યાખ્યાયિત કરવા દે છે.
દાખલા તરીકે, $h(x) = \Add(x, x)$ને આ રીતે વ્યાખ્યાયિત કરી શકાય:
\[
h(x_0) = \Add(\Proj{1}{0}(x_0),\Proj{1}{0}(x_0)).
\]
અહીં $k=2$, $n=1$ છે. $f(y_0,y_1)$ની ભૂમિકા $\Add$ ભજવે છે;
$g_0(x_0)$ અને $g_1(x_0)$ બંનેની ભૂમિકા~$\Proj{1}{0}(x_0)$,
એટલે કે એક-આર્ગ્યુમેન્ટવાળું પ્રક્ષેપ વિધેય (તદેવ વિધેય), ભજવે છે.

જો $f(y_0, y_1)$ વિધેય પહેલેથી આપણી પાસે હોય, તો
$h(x_0, x_1) = f(x_1, x_0)$ને આ રીતે વ્યાખ્યાયિત કરી શકાય:
\[
h(x_0, x_1) = f(\Proj{2}{1}(x_0, x_1),\Proj{2}{0}(x_0, x_1)).
\]
અહીં $k=2$, $n = 2$ છે, અને $g_0$ તથા $g_1$ની ભૂમિકા અનુક્રમે
$\Proj{2}{1}$ અને~$\Proj{2}{0}$ ભજવે છે.

$g_0$, \dots,~$g_{k-1}$ બધાંને એકસરખા $n$ આર્ગ્યુમેન્ટ હોવા
જ જોઈએ, એ શરત અંગે પણ શંકા થઈ શકે. (યાદ રાખો કે વિધેયની
\emph{આર્ગ્યુમેન્ટ-સંખ્યા} એટલે તેના આર્ગ્યુમેન્ટોની સંખ્યા;
$n$-આર્ગ્યુમેન્ટવાળાં વિધેયની આ સંખ્યા~$n$ છે.) પરંતુ પ્રક્ષેપ
વિધેયો ઉમેરવાથી જરૂરી લવચીકતા મળે છે. દાખલા તરીકે, ધારો કે $f$
અને~$g$ $3$-આર્ગ્યુમેન્ટવાળાં વિધેયો છે, અને $h$ નીચે પ્રમાણે
વ્યાખ્યાયિત $2$-આર્ગ્યુમેન્ટવાળું વિધેય છે:
\[
h(x,y) = f(x,g(x,x,y),y).
\]
પ્રક્ષેપ વિધેયો વડે $h$ની વ્યાખ્યા આ રીતે ફરી લખી શકાય:
\[
h(x,y) = f(\Proj{2}{0}(x,y), g(\Proj{2}{0}(x,y), \Proj{2}{0}(x,y),
\Proj{2}{1}(x,y)), \Proj{2}{1}(x,y)).
\]
આથી $h$, $f$નું $\Proj{2}{0}$, $l$ અને $\Proj{2}{1}$ સાથેનું
સંયોજન છે, જ્યાં
\[
l(x,y) = g(\Proj{2}{0}(x,y),\Proj{2}{0}(x,y),\Proj{2}{1}(x,y)),
\]
એટલે કે $l$, $g$નું $\Proj{2}{0}$, $\Proj{2}{0}$ અને
~$\Proj{2}{1}$ સાથેનું સંયોજન છે.

\end{document}
